\documentclass[11pt]{article}
\usepackage[margin=1in]{geometry}
\usepackage{amsmath,amssymb,amsthm,hyperref}
\begin{document}
\begin{center}{\Large A nonzero class above the predicted top degree}\\[4pt]
Disproof of conjecture 00000001649\\[4pt]
ziangni-sys --- October 8, 2026\end{center}

Conjecture 00000001649 predicts that the largest degree of nonzero rational
cohomology of $SL_n(\mathbb Z)$ is
\[
 t(n)=\left\lfloor\frac{n^2-1}{4}\right\rfloor-1.
\]
There is no exclusion of $n=2$ in either language version. For this standard
case, integer arithmetic gives
\[
 t(2)=\left\lfloor\frac34\right\rfloor-1=-1.
\]
We exhibit a genuine nonzero cohomology class in degree zero, which lies
strictly above this predicted top degree.

\paragraph{The actual group and its cochains.}
Take
$G=SL_2(\mathbb Z)=\{M\in M_2(\mathbb Z):\det M=1\}$, with ordinary matrix
multiplication, and trivial rational coefficients. In its inhomogeneous
group cochain complex, $C^0=\mathbb Q$, $C^1=\operatorname{Map}(G,\mathbb Q)$,
and the differential is
\[
 (d^0q)(h)=h\cdot q-q=q-q=0.
\]
Thus $\ker d^0=\mathbb Q$. Since there are no negative-degree cochains,
$\operatorname{im}d^{-1}=0$, so
\[
 H^0(SL_2(\mathbb Z);\mathbb Q)
 =\ker d^0/\operatorname{im}d^{-1}
 =\mathbb Q/0\cong\mathbb Q.
\]
The class of the constant zero-cochain $1$ is nonzero: its image under this
isomorphism is $1\ne0$. In particular, $H^0$ has dimension one.

\paragraph{Contradiction.}
By the conjecture's definition, a largest nonzero degree $d$ must satisfy
$i\le d$ whenever $H^i(SL_2(\mathbb Z);\mathbb Q)\ne0$. The class above
therefore forces $0\le d$. The proposed value $d=t(2)=-1$ instead forces
$0\le-1$, a contradiction. This argument needs no calculation of higher
cohomology and no assumption about its vanishing. A single nonzero class
above a proposed maximum refutes that maximum.

\paragraph{Scope and arithmetic.}
The subtraction in the displayed conjecture is ordinary integer
subtraction; $\lfloor3/4\rfloor-1=-1$, rather than natural-number truncated
subtraction. The asserted formula is already false at $n=2$, independently
of the additional virtual-cohomological-dimension and boundary descriptions.
An amended formula restricted to larger $n$ would be a different statement.

\paragraph{Lean certificate.}
The project uses Mathlib's actual $2\times2$ integer special linear group,
constructs trivial-action group cochains, the degree-zero differential,
cocycles and the zero-boundary quotient, and gives an explicit linear
equivalence with $\mathbb Q$. It certifies a nonzero class, dimension one,
the integer floor computation and the impossibility of any proposed top
degree equal to $t(2)$ satisfying even its necessary degree-zero bound.
\end{document}
